\chapter{Conclusion et perspectives}
\label{chap:conclusion}

\section{Synthèse des contributions}

Ce projet a abouti à un interpréteur complet et fonctionnel pour le mini-langage
impératif \ilang. L'ensemble de la chaîne demandée a été réalisé : analyse
lexicale par Flex, analyse syntaxique LALR(1) par Bison avec construction d'un
arbre syntaxique abstrait, table des symboles gérant la portée par une pile
d'environnements, deux passes d'optimisation, et un évaluateur récursif gérant
variables, expressions, structures de contrôle, fonctions récursives et
entrées-sorties. La gestion des erreurs couvre toutes les catégories du cahier des
charges, chacune localisée à la ligne.

À cette base se sont ajoutées trois extensions, à savoir les nombres flottants,
les instructions de rupture de boucle \code{arrêter} et \code{continuer}, et la
localisation systématique des erreurs. S'y est jointe une exploration assumée
hors du périmètre demandé, un second back-end générant du code assembleur i386,
dont les sorties coïncident exactement avec celles de l'interpréteur. La qualité d'ingénierie visée a été tenue : zéro
avertissement, zéro conflit grammatical, zéro fuite mémoire.

\section{Prolongements possibles}

Plusieurs directions permettraient d'enrichir le projet, et l'architecture
retenue rend la plupart d'entre elles abordables sans refonte. Je les présente en
indiquant à chaque fois ce qu'elles impliqueraient concrètement, car c'est cette
estimation du coût qui distingue une vraie perspective d'un simple v\oe u.

\paragraph{Chaînes de caractères.} Ajouter un type \code{string} serait sans doute
l'extension la plus naturelle après les flottants, car elle suit exactement le
même schéma. Il suffirait d'enrichir l'union étiquetée \code{Value} d'une
troisième variante, et de gérer l'allocation et la libération des chaînes :

\begin{lstlisting}[style=c,caption={Esquisse d'extension de \code{Value} aux chaînes.}]
typedef enum { VAL_INT, VAL_FLOAT, VAL_STR } ValType;
typedef struct Value {
    ValType type;
    long    ival;
    double  fval;
    char   *sval;   /* nouveau : pointeur sur la chaine (a liberer) */
} Value;
\end{lstlisting}

La vraie difficulté ne serait pas le type lui-même mais la gestion de la durée de
vie des chaînes, qui demanderait soit un comptage de références, soit un
ramasse-miettes, pour éviter fuites et libérations prématurées.

\paragraph{Tableaux.} Introduire des tableaux, avec une syntaxe du type
\code{tab = [1, 2, 3]} et un accès indexé \code{tab[i]}, demanderait deux nouveaux
types de n\oe uds dans l'AST, l'un pour le littéral de tableau et l'autre pour
l'indexation, ainsi qu'une nouvelle variante de \code{Value} pointant vers un bloc
alloué. L'évaluateur s'enrichirait de deux cas dans son \code{switch}, ce qui
reste dans la logique du parcours d'arbre existant.

\paragraph{Machine virtuelle à bytecode.} Pour dépasser les deux principales
limites de l'interpréteur, la lenteur du parcours d'arbre et la profondeur de
récursion bornée, la voie classique consiste à compiler l'AST vers un jeu
d'instructions de bas niveau, exécuté par une boucle dédiée disposant de sa propre
pile explicite. On définirait d'abord un jeu d'instructions :

\begin{lstlisting}[style=c,caption={Esquisse d'un jeu d'instructions de bytecode.}]
typedef enum {
    OP_PUSH, OP_LOAD, OP_STORE,      /* constantes et variables */
    OP_ADD, OP_SUB, OP_MUL, OP_DIV,  /* arithmetique sur la pile */
    OP_JMP, OP_JZ,                   /* sauts (controle de flot) */
    OP_CALL, OP_RET, OP_PRINT        /* fonctions et affichage */
} ByteOp;
\end{lstlisting}

L'évaluation reviendrait alors à une boucle lisant ces instructions une à une.
C'est l'approche de Python ou de Lua, et elle supprimerait du même coup la limite
de récursion, puisque la pile d'appels deviendrait une structure de données gérée
explicitement plutôt que la pile du langage hôte.

\paragraph{Back-end assembleur complet, vérification de types, afficheur d'AST.}
Trois autres prolongements me semblent intéressants. Le back-end assembleur
pourrait être étendu aux flottants, au prix d'une gestion du jeu d'instructions
flottantes du processeur, et reproduire la gestion « douce » de la division par
zéro en émettant un test avant chaque division. Une passe d'analyse sémantique
préalable pourrait signaler davantage d'erreurs avant même l'exécution, par
exemple l'usage d'une variable jamais affectée. Enfin, un afficheur d'AST
reconstituant le source avec un parenthésage minimal serait un outil précieux,
aussi bien pour le débogage que pour l'enseignement, car il rendrait visible la
structure que l'analyseur a réellement construite.

\section{Applications potentielles}

Au-delà de l'exercice, les techniques mises en \oe uvre dépassent largement la
construction de langages généralistes. Elles s'appliquent à tout traitement de
texte structuré : langages de configuration, langages dédiés (\emph{DSL}) métier,
moteurs de calcul de formules (tableurs), analyseurs de protocoles, ou encore
préprocesseurs et transpileurs. La compétence centrale, qui consiste à décrire
formellement une syntaxe puis à lui associer une sémantique, est directement
réutilisable dans tous ces contextes, bien au-delà du seul cas des langages de
programmation généralistes.

\section{Recommandations}

Pour qui reprendrait ou referait un tel projet, trois recommandations se dégagent
de l'expérience. Premièrement, \textbf{construire incrémentalement} : un c\oe ur
minimal qui fonctionne de bout en bout vaut mieux qu'un édifice complet jamais
exécuté. Deuxièmement, \textbf{tester tôt et systématiquement}, y compris les cas
limites et la mémoire (Valgrind), plutôt que de se fier aux seuls exemples
heureux. Troisièmement, \textbf{exploiter les outils} plutôt que de les
contourner : laisser Bison signaler les conflits, laisser le compilateur signaler
les avertissements, et traiter chacun comme une information à comprendre.

\section{Apports personnels et conclusion}

Sur le plan des apprentissages, ce projet a transformé des notions jusque-là
théoriques en savoir-faire concret. La distinction entre niveaux régulier et
hors-contexte, la mécanique d'une grammaire LALR(1), le rôle d'un arbre syntaxique
comme représentation pivot, la gestion d'une pile d'appels : autant de concepts
qui prennent un sens entièrement nouveau une fois qu'on les a implémentés et
débogués soi-même, plutôt que seulement lus dans un cours. Le moment le plus marquant aura été de voir le même arbre produire,
selon le back-end choisi, soit un résultat calculé, soit des instructions
machine : cette dualité résume à elle seule ce que le cours cherche à transmettre.

En définitive, \ilang{} n'est pas un langage destiné à être utilisé, mais un
langage destiné à être \emph{compris}. C'est à cette aune qu'il faut juger le
projet : non à la richesse de ses fonctionnalités, mais à la clarté avec laquelle
il rend visible le chemin qui mène d'un texte source à son exécution.
